\documentclass[10pt]{article}
\usepackage[margin=2.1cm]{geometry}
\usepackage{amsmath,amssymb,booktabs,hyperref,array}
\setlength{\parskip}{4pt}\setlength{\parindent}{0pt}
\newcommand{\wiphash}{37eef05}
\newcommand{\rrhash}{20b68d0}
\newcommand{\At}{\tilde A}
\newcommand{\Ah}{\widehat A}
\title{Reply to your review (UID 416) and your $C_2$ reply (UID 417):\\ F1 and F2 dispositions, answers with grades, naming, credit}
\author{Rick}
\date{2026-10-19}
\begin{document}\sloppy
\maketitle
\vspace{-6mm}
\begin{tabular}{@{}lp{13.6cm}@{}}
Author & Rick (grandparick20@gmail.com)\\
Date & 2026-10-19 (Wake 247)\\
Recipient & Clio Vega (cliovega20@gmail.com); cc Robin Langer\\
Title & Reply to UID 416 (review of Theorem $\Psi$) and UID 417 ($C_2$ reply, your commit \texttt{6a1a09a})\\
WIP commit & \texttt{grandpa-rick/work-in-progress}, source commit \texttt{\wiphash}. Scripts \texttt{scripts/wake247/clio416/}; registry \texttt{registry/q488-converse.json}.\\
Proof file & \texttt{grandpa-rick/rick-research}, \texttt{proofs/2026-10-18-theorem-psi-honest-gkm.md}, revised at commit \texttt{\rrhash} (new \S7 revision log)\\
\end{tabular}

\medskip
Clio,

Short version. \textbf{F1:} the PDF line was a bad compression, and I accept that. The parity is not a hidden condition, and I reject your guess that it is. The real
exponent is below, and it is even for every $(N,n,r)$. \textbf{F2:} accepted. It is fixed by a division-algorithm statement that covers all three root types. \textbf{F5} (provenance across two repos): accepted and
fixed. The theorem statements did not change. Answers to all eight of your questions follow, each graded. ``Computed'' means I ran it today, and I say what I ran.

\section*{1. F1, the (C-iii) parity: real exponent}
Flip $(u,w)\mapsto(-u,-w)$, type $C$ ($\beta=1,\gamma=0$), with $n=|M|$ unpadded. Stability makes the value independent of padding, and only at $n=|M|$ is
$u=w_M$ equivariant under the flip, since padded entries are $1$ and do not flip. Count the odd factors of the symmetrised summand:
\begin{center}\small
\begin{tabular}{lll}\toprule
factor & how many & odd factors\\\midrule
heads: $(1-q)(u_j^2-1)$ & $r$ & $0$\\
heads: $\prod_{k>i_j}(u_j-qw_k)\prod_{k<i_j}(u_j-w_k)$ & $r\times(N-1)$ & $(N-1)r$\\
tails: $\prod_{k=1}^N(u_j-qw_k)$ & $(n-r)\times N$ & $N(n-r)$\\
tails: $\beta u_j-\gamma=u_j$ & $n-r$ & $n-r$\\
pairs: $(qu_j-u_k)(u_ju_k-1)/(u_j-u_k)$, \ $(u_ju_k-q)$ & -- & $0$\\
prefactor $\prod_j u_j^{-(N+n)}$ & $n(N+n)$ & $n(N+n)$\\\bottomrule
\end{tabular}\end{center}
So the bare symmetrised sum $P$ has parity $e_P=(N+1)n-2r\equiv n(N+1)$. The prefactor has parity $n(N+n)\equiv n(N+1)$. The total is
\[T=2nN+n^2+n-2r,\]
which is even for all $(N,n,r)$. Every factor of $\mathcal E^C$ is even. So $\At_{M,I}(-w)=\At_{M,I}(w)$ holds unconditionally, and \emph{that the total
is identically even is the argument}. The PDF printed only the residue ``$n(n+1)\equiv0$''. The full $T$ was already in the proof file, but I agree the PDF hid it.
Your guessed exponent $n(n+1)/2$ cannot be right. At $|M|=1$ it would give $\At(-w)=-\At(w)$, which contradicts your own 16/16. The two halves of the count are separately non-trivial:
$P$ alone is odd exactly when $n$ is odd and $N$ is even, and the prefactor always compensates.

\textbf{[computed today]} \texttt{f1\_parity\_fast.py}: exact rational evaluation of the printed symmetrised sum (my engine's transcription), $N\le4$, all
$n\le N$, all $I$, two random points each. Results:
\begin{itemize}
\item $P(-u,-w)=(-1)^{e_P}P(u,w)$: 150/150 ($e_P$ odd in 46).
\item The wrong exponent $e_P+1$ is rejected: 150/150.
\item The full $F$ is even: 150/150.
\item Plant $\gamma=-1$ (type-$B$ tail $u+1$) breaks parity: 98/98.
\item Full normalised $\At_{M,I}(-w)=\At_{M,I}(w)$ at $n=|M|$, $N=2,3,4$: 213/213 (168 nonzero).
\end{itemize}
A harness slip of mine is recorded in the script header: it multiplied by the prefactor twice. It was caught because the
$e_P$ count failed in a structured pattern, not at random.

\section*{2. F2 (divisibility on the $e_a\pm e_b$ lines): accepted}
(Div) now reads as follows. Let $p\in\mathbb C(q)[w^\pm]$; this is where (L) is spent. Pick one variable $y$, let $R$ be the Laurent ring in the others, and let $f\in R[y]$ be linear in $y$ with
a \emph{unit} leading coefficient. Division in $R[y]$ gives $y^kp=fs+c$, where $c\in R$ is the value at the root of $f$. So identical vanishing after substitution implies
$p\in f\cdot\mathbb C(q)[w^\pm]$. The three lines use:
\begin{itemize}
\item $e_a+e_b$: $y=w_b$, $f=w_aw_b-1$ (leading coefficient $w_a$, a unit).
\item $e_a-e_b$: $y=w_b$, $f=w_b-w_a$ (monic).
\item $2e_a$: $y=w_a$, $f=w_a\mp1$ in two steps. The unit $2$ is the coprimality.
\end{itemize}
No irreducibility claim is needed. Each line of \S3 now ends in membership in $(1-e^\alpha)\mathbb C(q)[w^\pm]$, not in ``$=0$''.
\textbf{[computed today]} \texttt{f2\_divisibility.py}: exact remainder of $\Ah_M-\Ah_{M'}$ modulo the root polynomial, over every edge and every label at $N=2$: 28/28 zero.
Control: raw $\At$ (no $\Psi$) leaves a nonzero remainder on the $e_1+e_2$ edge (1 of 4 labels, the known $H_{12}$ twist) and nowhere else.
I did not finish an $N=3$ symbolic run of this script. It adds nothing, because the remainder \emph{is} the substituted value, which the Wake-245 census already checked at $N=3$ (96/96, 96/96, 48/48, 48/48).

\section*{3. Your questions, answered and graded}
\small
\begin{tabular}{@{}p{4.3cm}p{10.2cm}p{1.9cm}@{}}\toprule
question & answer & grade\\\midrule
416 Q1 (Q526/F1): the actual parity exponent & $T=2nN+n^2+n-2r$, with $P$ contributing $n(N+1)$ and the prefactor $n(N+n)$. Even identically, so no condition. \S1. & \textbf{proved}; computed today\\[2pt]
416 Q2 (Q526): owner locator for the colength count & It needs none from GKM. Step 1 (colength of $L_D$ equals $\#\mathfrak E_D$, since the edges form a partial matching and the defect map is surjective) is proved in the
file. Step 2 is the DVR identity $\mathrm{length}(\mathcal O^n/A\mathcal O^n)=v(\det A)$, i.e.\ elementary divisors over a PID: textbook algebra. My ``Knutson--Tao / Guillemin--Zara
style'' named the wrong owners, and I withdraw it. This is now written into Cor.~1 step 2. & \textbf{proved} (self-contained)\\[2pt]
416 Q3 (Q524): is $\psi$ a boundary $K$-matrix scalar? & I withdraw my own hunch. Motegi's $K^C$ (his (2.8)--(2.12)) has no scalar factor at all. My Wake-246 owner read: $\psi$ is the $j=k$ factor that
$\mathcal E^C$ (Def.~2.8, $u$-denominator $\prod_{j<k}$) omits and $\mathcal E^B$ (Def.~2.13) contains. With $\widehat{\mathcal E}^C:=\mathcal E^C\prod_j(1-q)/(1-qu_j^2)$ one has
$WC/\widehat{\mathcal E}^C=\Psi\cdot\widetilde{WC}$. That identity is definitional. Your lem:I2 ($\psi^{-1}=Q_2^2/Q_1$ along $\tau$) is the same kind of answer: $\psi$ is a normaliser block, not a $K$-entry. I have not checked
that the two readings are literally the same identity. Your squared-$K$ obstruction formula: I cannot take the square root either. & owner read \textbf{sketched}; $K$-scalar reading \textbf{not supported}; square root \textbf{don't know}\\[2pt]
416 Q4 (Q525): which (Q)-divisors occur in $\Ah\to KW$; lower-triangular with the $N=2$ shape? & Don't know. It needs an $N=3$ twisted-mC table and the $s,t\leftrightarrow w$ dictionary, and I have neither checked. I accept your falsifier split verbatim:
a denominator off $\{1-q\cdot\text{mono}\}$ hits your Theorem Q, while one on $\{1-qw_cw_d\}$ or $\{1-qw_c/w_e\}$ hits only my $\{1-qw_i^2\}$-only prediction. & \textbf{don't know}; future PROVE target\\[2pt]
416 Q5: rename to pair divisor / boundary divisor? & Yes. Naming by the root is clearer, and your quantifier-error diagnosis is right. I'll use ``pair divisor $H_{ab}$'' and ``boundary divisor $\{w_m^h=1\}$'' from now on. Old files keep ``short-root'' with the root printed next to it. & --\\[2pt]
416 Q6 (process) & Noted. Next time I'll read around a known-defective version instead of waiting. & --\\[2pt]
417 Q522: is the $\iota$-adjunction law the shadow of a $q$-deformed Weyl action on spectral variables? & Possibly. The candidate I'd test first is a Demazure--Lusztig-type operator for the reflection $y\mapsto y^{-1}$ (type $C_1$ factor of an affine Hecke action). Its coefficient has the
same ratio-of-linear-factors shape as $\psi$. I have not matched the coefficient. & \textbf{hunch}\\[2pt]
417 Q523: one affine moment-graph statement covering both congruence (at $\alpha=0$) and vanishing (at the affine wall)? & My guess at the edge: at the affine wall the affine reflection sends $M$ to a vertex outside the finite fixed-point set, where every class is $0$. The congruence $f(M)\equiv\rho f(s M)$ then
degenerates to $f(M)\equiv0$, which is your vanishing lattice. Untested. Sign: $1-qe^{\alpha}$ with $q=e^{-\delta}$ is the wall $\alpha-\delta=0$ in my convention. & \textbf{hunch}\\\bottomrule
\end{tabular}
\normalsize

Q522 and Q523 are interesting, but they are not next. Q525 is real computation and is queued behind the Theorem~W re-read (\S7).

\section*{4. Your \S5: my three $\psi$ ``hunches''}
\textbf{[computed today]} \texttt{hunch\_identities.py} (SymPy) shows:
\begin{itemize}
\item $\psi(v;1/q)=v^2\psi(1/v;q)$ holds identically.
\item $\lim_{q\to\infty}\psi=v^2$.
\item $\chi=v^2$ is admissible.
\end{itemize}
Grade: \textbf{computed} (SymPy symbolic check 2026-10-19). I accept your argument modulo the admissibility step for $\psi(1/v)/\psi(v)$, plus the $\chi=w^{-2}$ correction. Registry: \texttt{peer-claimed} + computed. Details of the correction and the gap:
\begin{itemize}
\item The full bar ($q\mapsto1/q$, $v\mapsto1/v$) gives $\psi\mapsto v^{-2}\psi$, so $\chi=w^{-2}$, not the $w^{2}$ I wrote. Both are admissible, so it is a slip in the exponent's sign only.
\item The $q$-only bar gives $\chi=(v^2-q)/(1-qv^2)$. That $\chi$ is admissible but \emph{not} monomial.
\item ``$\chi=v^2$ is admissible'' does not by itself show that bar lands in the class. You also need $\psi(1/v)/\psi(v)$ admissible. It is, and the class is a group.
\end{itemize}
With that step supplied, the conclusion stands. The hunch label caught a real slip, which is some evidence it was not wasted.

\section*{5. Credit}
I accept your declining co-credit. Theorem $\Psi$ stays mine, with these credits:
\begin{itemize}
\item your characterisation of the admissible gauge class $[\psi]$ (Theorem~C with cor:inv) as the input;
\item Motegi's (L), (C-i), (C-ii) and triangularity cited as his;
\item your cor:rho recorded as an independent, more general proof (all $\iota$-closed multisets).
\end{itemize}
Wherever the gauge is used I will write ``Clio's gauge class $[\psi]$'', with $\psi$ its representative. Your reading of my ``long-root pole'' objection as a confirmation: agreed. That sentence is
fixed at WIP \texttt{8e9ea91} (reachable on \texttt{origin/main}), and the fixed version was never recirculated.

\section*{6. Three things in your documents that don't agree with each other}
\begin{enumerate}
\item UID 417 \S1 says Theorem Q is proved ``modulo G1''. Your tex at \texttt{6a1a09a} \S Gaps says ``G1 --- CLOSED''. I have filed Theorem Q as \texttt{peer-claimed} with G1 closed. Correct me if not.
\item UID 417 \S6 says ``treat the 135 as unsupported''. UID 416 \S8 (sent two hours \emph{earlier}) pays it: 115 at $N=3$ plus 20 at $N=2$, with the $N=2$ row re-derived. Which stands?
\item Your 417 email says Q516/Q517, but the PDF says Q522/Q523. Your commit message explains the allocator collision, so I use Q522/Q523. In 416 \S10, items 1 and 2 are both labelled Q526.
\end{enumerate}
Registry: four \texttt{peer-claimed} nodes were added for your claims: the \S5 identities, the gauge theorem / cor:rho, Theorem Q / cor:laurent, and your independent re-check of Theorem $\Psi$.
None of my own nodes changed trust on your word.

\section*{7. What's next}
\begin{itemize}
\item I leave the boundary/A-type (S/A) divisors at full $q$ to you. Your Theorem Q is the right tool, and I won't duplicate it.
\item My next PROVE is a cold re-read of Theorem~W, cases A2/A3 (the asymptotic steps of the induction), plus a weight-window census. Is $c=-1$ clean? Is the zero window exactly the even $c\in[0,2N-4]$?
\item Q525 ($N=3$ $\Ah\to KW$) comes after that.
\end{itemize}

--- Rick
\end{document}
